%% =============================================================================
%% Paper 1 - ALL FIGURES (for main_final.tex, EDHEC format)
%% =============================================================================
%% Figures placed after Tables, on single pages. Captions at bottom.
%% =============================================================================

%% -----------------------------------------------------------------------------
%% SECTION 3 - Data, Strategy Construction, and Performance
%% -----------------------------------------------------------------------------

%% Figure: TIPS--Treasury basis, pre-2009 vs. today (referenced in sec3 BTP-Italia
%% paragraph; two side-by-side panels in a single float, JoF-style).
\clearpage
\begin{figure}[p]
\centering
\begin{minipage}[t]{0.48\textwidth}
\centering
\includegraphics[width=\textwidth, height=5.5cm, keepaspectratio=false]{\figpath/TIPSpre2009.pdf}
\caption[TIPS--Treasury basis, 2004--2009]{TIPS--Treasury basis, 2004--2009 (from \citealp{fleckenstein2014tips}).}
\label{fig:tips_pre2009}
\end{minipage}
\hfill
\begin{minipage}[t]{0.48\textwidth}
\centering
\includegraphics[width=\textwidth, height=5.5cm, keepaspectratio=false]{\figpath/TIPS_today.pdf}
\caption[TIPS--Treasury Basis, 2009--Present]{TIPS--Treasury basis, 2009--present. The mispricing has vanished.}
\label{fig:tips_today}
\end{minipage}
\end{figure}

%% Figure: raw market-level basis time series (one panel per strategy)
\clearpage
\begin{figure}[p]
\centering
\includegraphics[width=\textwidth]{\figpath/basis_time_series.pdf}
\caption[Market-Level Basis Time Series]{Market-Level Basis Time Series for Fixed Income Arbitrage Strategies.
For BTP~Italia, the basis is the cross-sectional mean of all outstanding
issues with at least six months to maturity. For iTraxx~Combined, it is
the mean of the four on-the-run sub-index skews. For CDS--Bond~Basis,
it is the mean of the $N=20$ most negative single-name bases in the
iTraxx~Main universe on each day (only bonds with negative basis are
included). Dashed lines indicate entry thresholds.}
\label{fig:basis_time_series}
\end{figure}

%% Figure: cumulative return indices and drawdowns
\clearpage
\begin{figure}[p]
\centering
\includegraphics[width=\textwidth]{\figpath/equity_curves_combined.pdf}
\caption[Cumulative Return Indices and Drawdowns]{Cumulative Return Indices and Drawdowns of Fixed Income Arbitrage
Strategies (Signal-Weighted). Each row shows one strategy: the upper
sub-panel reports the cumulative return index from inception, and the
lower sub-panel reports the running drawdown (peak-to-trough). Sample
periods differ across strategies: BTP~Italia from 2014, iTraxx Combined
from 2008, CDS--Bond Basis from 2005. All series end in May~2025 and
are net of transaction costs.}
\label{fig:equity_curves}
\end{figure}

%% -----------------------------------------------------------------------------
%% SECTION 4 - Benchmark Factor Models
%% -----------------------------------------------------------------------------

%% Figure: Rolling 36-month alpha across benchmark factor models
\clearpage
\begin{figure}[p]
\centering
\includegraphics[width=\textwidth]{\figpath/rolling_alpha_composite_monthly.pdf}
\caption[Rolling Alpha Across Benchmark Factor Models]{Rolling 36-Month Alpha Across Benchmark Factor Models.
  Each panel shows one strategy. Lines represent rolling alpha under
  Duarte et al.\ (2007, blue), Brooks et al.\ (2020, orange), and
  Fung \& Hsieh (2004, green). Red shading: HIGH stress regime
  (iTraxx Main 5Y $> 100$ bps).}
\label{fig:rolling_alpha_composite}
\end{figure}

%% -----------------------------------------------------------------------------
%% SECTION 5 - Factor Selection and Spanning Regressions
%% -----------------------------------------------------------------------------

%% Figure: PCA scree plot — variance explained by principal components
\clearpage
\begin{figure}[p]
\centering
\includegraphics[width=0.85\textwidth]{\figpath/pca/pca_scree_plot.pdf}
\caption[PCA Scree Plot]{Scree Plot: Cumulative Variance Explained by Principal Components.
  Dashed line indicates 80\% threshold. Baseline specification uses
  $K=8$ components.}
\label{fig:pca_scree}
\end{figure}

%% -----------------------------------------------------------------------------
%% SECTION 6 - Cross-Strategy Interdependencies
%% -----------------------------------------------------------------------------

%% Figure: Pairwise correlations by stress regime (bar chart)
\clearpage
\begin{figure}[p]
\centering
\includegraphics[width=0.9\textwidth]{../results/rq3_duffie/figures/C1_regime_correlations_bar.pdf}
\caption[Pairwise Correlations by Stress Regime]{Pairwise Correlations by Stress Regime.
  Bars show Pearson correlations of $\Delta m$ in LOW, MEDIUM, and HIGH
  stress regimes defined by iTraxx Main 5Y thresholds (60/120 bps).}
\label{fig:rq3_regime_bar}
\end{figure}

%% Figure: Generalized Impulse Response Functions (Pesaran-Shin 1998)
\clearpage
\begin{figure}[p]
\centering
\includegraphics[width=\textwidth]{../results/rq3_duffie/figures/E1_girf.pdf}
\caption[Generalized Impulse Response Functions]{Generalized Impulse Response Functions.
  Responses of $\Delta m$ to one-standard-deviation shocks.
  Shaded areas: 95\% bootstrap confidence intervals (1000 replications).}
\label{fig:rq3_girf}
\end{figure}